%
% File: abstract.tex
% Author: V?ctor Bre?a-Medina
% Description: Contains the text for thesis abstract
%
% UoB guidelines:
%
% Each copy must include an abstract or summary of the dissertation in not
% more than 300 words, on one side of A4, which should be single-spaced in a
% font size in the range 10 to 12. If the dissertation is in a language other
% than English, an abstract in that language and an abstract in English must
% be included.

\chapter*{Abstract}
\initial{T}oday's enterprises handle a wide variety of documents, including contracts, invoices, payslips and CVs. Extracting data from such heterogeneous documents is a major challenge for information systems. Existing solutions are often inadequate: rule-based approaches are rigid, while artificial intelligence models frequently sacrifice accuracy on numerical values and positional context. To address these limitations, this research introduces an innovative hybrid methodology that combines template-based extraction, named entity recognition (NER), and pattern recognition techniques. Within an agile framework, the proposed system incorporates a continuous learning mechanism that significantly improves extraction accuracy over successive business iterations, adapting to new document formats without requiring extensive retraining. The prototype, integrated as a module into DOLibaar, reduced manual data entry by 80\% and achieved over 95\% extraction accuracy on a test set of 10,000 documents. In comparison, a regex-only approach plateaus at 72\% accuracy, while a generic, non-fine-tuned LLM reaches only 81\% on the same dataset, demonstrating the superiority of the hybrid method. The solution integrates into business systems through standard formats such as JSON and XML, with an average processing time of less than 500~ms per file. Although handling low-quality scanned documents remains an identified limitation, the modular architecture paves the way for adding OCR preprocessing, thereby supporting the industrialization of the solution.
\paragraph{Keywords: Automatic data extraction, Administrative documents, Named Entity Recognition (NER), Hybrid methodology, Continuous learning, DOLibaar.}

\clearpage
